\chapter{Las tres condiciones de derivabilidad}
\label{cap:d123}

El primer teorema necesita que $\Prov$ exista y diga la verdad en una dirección. El segundo
necesita mucho más: que la teoría sepa \emph{razonar} sobre su propia demostrabilidad. Esas
exigencias son tres, se llaman D1, D2 y D3, y no las publicó Gödel — son el artículo que anunció y
nunca escribió, y las suministraron Hilbert y Bernays en 1939.

Este capítulo dice qué pide cada una, cuáles son teoremas aquí y por qué la tercera lleva meses
resistiéndose. Y da una respuesta concreta a la pregunta que el lector debería estar haciéndose:
\emph{¿cuánto falta, exactamente?}

\section{Las tres, en una página}

\begin{matematica}[Lo que pide cada una]
\[
\begin{array}{ll}
\text{D1} & \vdash \varphi \;\Longrightarrow\; \derives \Prov(\codigo{\varphi}) \\[4pt]
\text{D2} & \derives \Prov(\codigo{\varphi \Rightarrow \psi}) \Rightarrow
             \bigl(\Prov(\codigo{\varphi}) \Rightarrow \Prov(\codigo{\psi})\bigr) \\[4pt]
\text{D3} & \derives \Prov(\codigo{\varphi}) \Rightarrow \Prov(\codigo{\Prov(\codigo{\varphi})})
\end{array}
\]

\noindent D1: \emph{si algo se demuestra, la teoría lo sabe.} \quad
D2: \emph{la demostrabilidad respeta el modus ponens.} \quad
D3: \emph{la teoría sabe que sabe.}
\end{matematica}

La diferencia de naturaleza entre la primera y las otras dos es la que gobierna todo el capítulo, y
conviene fijarla antes de mirar ningún código.

\begin{leccion}[Externo contra provable]
\textbf{D1 es una afirmación de la metateoría}: para \emph{cada} $\varphi$ demostrable,
\emph{nosotros} construimos la derivación de $\Prov(\codigo{\varphi})$. Se demuestra caso a caso, y
cada caso es finito.

\textbf{D2 y D3 son afirmaciones de la teoría objeto}: la que se demuestra es una fórmula con un
$\Rightarrow$ dentro, cuantificada sobre \emph{todas} las derivaciones posibles. Y razonar sobre
todas las derivaciones, dentro de la teoría, exige inducción — que es justo lo que
Q\texttt{++} no tiene.

Ésa es toda la historia. D1 es trabajo; D3 es un problema.
\end{leccion}

\section{D1: un teorema, sin hipótesis}

Ya está demostrada, y el capítulo~\ref{cap:d1} contó cómo. Aquí sólo el recordatorio de que no
arrastra nada:

\leanfrag{repr_pos_prf}

Es la $\Sigma_1$-completitud \textbf{externa}: dada una demostración concreta, se construye el
testigo concreto. Ninguna inducción de la teoría objeto interviene.

\section{D2: también un teorema, y es el que suele postularse}

\leanfrag{d2_prf}

\selee{«la teoría demuestra que, si es demostrable $A \Rightarrow B$ y es demostrable $A$, entonces
es demostrable $B$»}

\begin{observacion}
Merece la pena subrayarlo porque es donde muchas presentaciones —y una versión anterior de este
mismo proyecto— se ahorran el trabajo: \textbf{D2 aquí es un teorema, no un postulado}. La
diferencia se ve en el inventario: el capítulo~\ref{cap:lean} decía que un \ident{axiom} es
una deuda, y ésta está pagada.

La prueba, por cierto, es instructiva: elimina el $\exists$ externo del predicado de
demostrabilidad —el testigo es la propia derivación— y después concatena las dos cadenas. Es el
\emph{modus ponens} de la teoría, escrito dentro de la teoría.
\end{observacion}

\section{D3: la que Gödel anunció, y estuvo postulada hasta anteayer}

Durante casi toda la vida de este proyecto, D3 fue lo que la mayoría de las presentaciones hacen
con ella: un postulado. Estaba escrita como \ident{axiom} de Lean —no como axioma objeto de la
lista de 141, sino como un supuesto de la \textbf{metateoría}—, lo que significaba que todo
teorema que la citara la llevaba impresa en su \emph{footprint}. El segundo teorema de Gödel se
publicaba con esa marca.

Ya no. Desde el 10 de septiembre de 2026 es un teorema, con la misma firma exacta que tenía el
axioma:

\leanfrag{d3_teorema}

\selee{«la teoría demuestra que, si $\varphi$ es demostrable, entonces es demostrable que
$\varphi$ es demostrable»}

\begin{leccion}[Retirar un axioma sólo puede fortalecer lo que había]
Nótese qué es lo que \emph{no} cambió: ni el enunciado de D3, ni el de \ident{goedel_second'}, ni
una línea de su demostración. Lo único que cambió es de dónde viene D3 — antes se suponía, ahora se
deriva.

Ésa es la propiedad que hace que este tipo de deuda sea \textbf{segura de contraer}: mientras el
postulado esté a la vista y con la firma que tendrá el teorema, pagarlo después no obliga a
reconstruir nada aguas abajo. Se sustituye la palabra \ident{axiom} por \ident{theorem} y todos
los \emph{footprints} de la cadena adelgazan solos.
\end{leccion}

El inventario de axiomas de Lean del proyecto pasó, ese día, de siete a \textbf{seis}. De los que
quedan, ninguno es una deuda de este tipo: son anclas de codificación y esquemas de inducción
declarados.

\section{Por qué costó tanto}

La demostración real de D3 es la $\Sigma_1$-completitud \textbf{provable}: que la teoría
demuestre, de toda derivación, que su verificador la acepta. Y ahí hay que razonar sobre un objeto
de longitud arbitraria \emph{desde dentro}.

\begin{leccion}[Dónde estaba la dificultad, en una frase]
D1 mira \emph{una} derivación y construye \emph{un} testigo. D3 tiene que decir algo sobre
\textbf{todas} las derivaciones, y decirlo la teoría, no nosotros. La primera es una función; la
segunda es una inducción — y la inducción de la teoría objeto es exactamente el recurso que
Q\texttt{++} no tiene y que \capfuturo{la inducción como precio} tendrá que comprar.
\end{leccion}

\section{La forma de la demostración}

Lo que hay debajo del teorema de una línea es una reducción de varios pisos, y su forma final es
ésta: D3 se sigue de reflejar \textbf{las dos mitades del cuerpo} del verificador.

\leanfrag{d3_prf_of_halves}

\selee{«si se reflejan las dos mitades del cuerpo —la buena formación de cada línea y sus
premisas—, D3 queda demostrada»}

Las dos reciben lo mismo: una cadena $q$ y un índice $i$ dentro de ella. Y de la cadena sale lo que
las dos necesitan:

\leanfrag{prf_lineWF_of_chainOk}

\begin{center}\small
\begin{tabularx}{\linewidth}{@{}l >{\raggedright\arraybackslash}X l@{}}
\toprule
mitad & qué refleja & se cerró \\
\midrule
(a) & que la línea $i$-ésima está bien formada & 2026-09-10 \\
(b) & que sus premisas están donde deben & 2026-09-10 \\
\bottomrule
\end{tabularx}
\end{center}

La mitad (a) es la que costó el episodio del capítulo~\ref{cap:muro-substfc}: los veintiún
reflectores del verificador, de los que siete estuvieron bloqueados cuarenta días. La mitad (b)
necesitó tres piezas más, y la primera de ellas resume bien de qué va todo esto:

\leanfrag{pcc_eval_premsOf}

\selee{«si la línea $t$ está bien formada, la teoría demuestra que el operador \emph{premisas de}
aplicado a su código da el código de sus premisas»}

Veintiuna ramas, una por cada tipo de línea que el verificador admite, porque \textbf{las premisas
sólo se pueden evaluar sabiendo de qué regla se trata}. Con ella, la cota; con la cota, el chasis
inductivo interior; y con las tres, el ensamblaje:

\leanfrag{d3_prf_real}

\begin{observacion}
Este libro llegó tarde por dos días. El capítulo que se escribió el 10 de septiembre por la mañana
decía —midiendo, y correctamente— que la mitad (b) estaba abierta en tres piezas, y señalaba además
un hueco de ensamblaje que ningún documento del proyecto había visto: la mitad (a) estaba probada y
\textbf{no la consumía nadie}.

Las tres piezas se cerraron esa misma tarde, en cuatro commits, y el hueco de ensamblaje lo cerró
el último de ellos. \ident{d3_prf_real} es exactamente el teorema que faltaba: la mitad (a) y la
mitad (b), enchufadas.

Se deja escrito porque es la clase de dato que un libro sobre un proyecto vivo debe dar y casi
ninguno da: \textbf{cuánto tiempo hubo entre saber qué faltaba y que dejara de faltar}. Aquí,
unas horas.
\end{observacion}

\section{Lo que las tres condiciones significan juntas}

\begin{leccion}[Por qué esto es el artículo que Gödel no escribió]
El primer teorema necesita D1 y nada más. El segundo necesita las tres, y ésa es la razón de que
Gödel anunciara una segunda parte y no la publicara: la demostración detallada exige justamente
esto, y «esto» resultó ser considerablemente más trabajo de lo que cabía en una nota.

Noventa y cinco años después, en una formalización mecánica, la proporción se repite con precisión
incómoda: D1 y D2 salieron en semanas, y D3 —la que Gödel dejó anunciada— costó meses y fue la
última en caer. Eso no dice nada malo del proyecto: es la medida honesta de lo que Hilbert y
Bernays tuvieron que suministrar cuando la segunda parte no llegó.

Con las tres demostradas, el segundo teorema deja de publicarse con una condición al margen. Qué
significa eso exactamente —y qué hipótesis \emph{sí} le quedan, que no son ninguna de las
tres— es el asunto de \capfuturo{Gödel II}.
\end{leccion}
